\chapter{Gastric Adenocarcinoma}
\index{Gastric Adenocarcinoma}
\label{sec:Gastric Adenocarcinoma}

\section{Overview}
Gastric adenocarcinoma is the fifth most common cancer worldwide and the fourth leading cause of cancer-related death. The two histological subtypes are intestinal (associated with \textit{H.~pylori}, chronic atrophic gastritis, intestinal metaplasia, and dietary factors) and diffuse (associated with CDH1 mutations, affecting younger patients with worse outlook). This condition accounts for a significant proportion of cancer-related morbidity and mortality worldwide. Early detection through routine screening and advances in treatment have improved survival rates in recent decades. This type of cancer arises from uncontrolled cell growth in affected tissues, with potential to invade nearby structures and spread to distant sites through the bloodstream or lymphatic system. The incidence varies by geographic region, age, and genetic background. Screening programs have improved early detection rates in many populations. Histological classification and molecular subtyping guide treatment decisions and provide prognostic information. Multidisciplinary care involving surgeons, medical oncologists, radiation oncologists, and pathologists is standard for optimal management.
\section{Causes}
This condition happens because of multifactorial carcinogenesis involving \textit{H.~pylori} infection, high-salt diet, smoking, family history, and pernicious anemia. Cancer develops through a multistep process involving genetic mutations that drive uncontrolled cell growth and proliferation. These mutations may be inherited or acquired through exposure to carcinogens, radiation, or random errors during cell division. Cancer develops through accumulation of genetic and epigenetic alterations that drive uncontrolled cell proliferation, evade growth suppression, resist cell death, and enable replicative immortality. Key molecular pathways involved include tumor suppressor genes (p53, RB, APC), oncogenes (KRAS, MYC, EGFR), and DNA repair mechanisms. Environmental carcinogens such as tobacco smoke, ultraviolet radiation, certain chemicals, and ionizing radiation can induce DNA damage. Chronic inflammation and persistent infections (HPV, hepatitis B/C, H. pylori) contribute to cancer development in some sites.
\section{Risk Factors}
Advancing age Tobacco use (active and secondhand smoke) Excessive alcohol consumption Family history of cancer and known genetic syndromes Obesity and physical inactivity Unhealthy diet (processed meats, low fiber) Exposure to carcinogens (asbestos, benzene, radiation) Chronic infections (HPV, hepatitis B/C, HIV) Immunosuppression Hormonal factors (early menarche, late menopause)

\begin{itemize}[noitemsep]
  \item Advancing age
  \item Tobacco use (active and secondhand smoke)
  \item Excessive alcohol consumption
  \item Family history of cancer and known genetic syndromes
  \item Obesity and physical inactivity
  \item Unhealthy diet (processed meats, low fiber)
  \item Exposure to carcinogens (asbestos, benzene, radiation)
  \item Chronic infections (HPV, hepatitis B/C, HIV)
  \item Immunosuppression
  \item Hormonal factors (early menarche, late menopause)
\end{itemize}
\section{Signs and Symptoms}
Persistent lump or mass in the affected area Unexplained weight loss and loss of appetite Chronic fatigue and weakness Persistent pain that worsens over time Changes in bowel or bladder habits Unexplained fever or night sweats Unusual bleeding or discharge Persistent cough or hoarseness Changes in skin appearance (new mole, changing wart) Difficulty swallowing or persistent indigestion

\begin{itemize}[noitemsep]
  \item Persistent lump or mass in the affected area
  \item Unexplained weight loss and loss of appetite
  \item Chronic fatigue and weakness
  \item Persistent pain that worsens over time
  \item Changes in bowel or bladder habits
  \item Unexplained fever or night sweats
  \item Unusual bleeding or discharge
  \item Persistent cough or hoarseness
  \item Changes in skin appearance (new mole, changing wart)
  \item Difficulty swallowing or persistent indigestion
\end{itemize}
\section{Progression and Stages}
Cancer staging follows the TNM system: Tumor (size and extent), Node (lymph node involvement), and Metastasis (spread to distant sites). Stage 0: Carcinoma in situ (abnormal cells confined to the original location). Stage I: Early-stage, small tumor confined to the organ of origin. Stage II: Locally advanced tumor with possible regional lymph node involvement. Stage III: More extensive local disease with significant lymph node involvement. Stage IV: Advanced disease with distant metastases to other organs. Higher stages generally indicate more advanced disease and poorer prognosis.

\begin{itemize}[noitemsep]
  \item Treatment focuses on using a thin camera tube mucosal removal or surgical gastrectomy with lymphadenectomy for early-stage disease
  \item Perioperative chemotherapy (FLOT regimen), targeted therapy (trastuzumab for HER2-positive disease), and immunotherapy (nivolumab for PD-L1-positive tumors) for advanced disease. Unfortunately, the outlook is often poor for advanced disease but favorable for early-stage disease with appropriate treatment.
\end{itemize}
\section{Complications}
Metastatic spread to vital organs (brain, liver, lungs, bones) Malignant effusions (pleural, peritoneal, pericardial) Superior vena cava syndrome (chest tumors) Spinal cord compression (spinal metastases) Pathological fractures (bone metastases) Tumor lysis syndrome (rapid cell death during treatment) Paraneoplastic syndromes (hormonal, neurologic, hematologic) Treatment-related complications (chemotherapy toxicity, radiation damage) Infections due to immunosuppression Venous thromboembolism

\begin{itemize}[noitemsep]
  \item Metastatic spread to vital organs (brain, liver, lungs, bones)
  \item Malignant effusions (pleural, peritoneal, pericardial)
  \item Superior vena cava syndrome (chest tumors)
  \item Spinal cord compression (spinal metastases)
  \item Pathological fractures (bone metastases)
  \item Tumor lysis syndrome (rapid cell death during treatment)
  \item Paraneoplastic syndromes (hormonal, neurologic, hematologic)
  \item Treatment-related complications (chemotherapy toxicity, radiation damage)
  \item Infections due to immunosuppression
  \item Venous thromboembolism
\end{itemize}
\section{Diagnosis}
Doctors diagnose this based on upper endoscopy with biopsy, with staging using TNM classification, CT imaging, and laparoscopy for peritoneal assessment. Diagnosis typically involves imaging studies (CT, MRI, PET scans) to identify the location and extent of disease. Biopsy with histopathological examination remains the gold standard for confirming the diagnosis and determining the specific type and grade. Physical examination including palpation of mass and lymph node assessment Complete blood count and comprehensive metabolic panel Tumor markers specific to cancer type (e.g., PSA, CA-125, CEA, AFP) Imaging: Ultrasound, CT scan, MRI, PET-CT for staging Biopsy: Core needle, excisional, or endoscopic biopsy for histopathology Immunohistochemistry to determine tumor subtype and receptor status Molecular and genetic testing (EGFR, KRAS, BRCA, MSI status) Sentinel lymph node biopsy for surgical staging Bone marrow biopsy for hematologic malignancies Genetic counseling and testing for hereditary cancer syndromes

\begin{itemize}[noitemsep]
  \item Physical examination including palpation of mass and lymph node assessment
  \item Complete blood count and comprehensive metabolic panel
  \item Tumor markers specific to cancer type (e.g., PSA, CA-125, CEA, AFP)
  \item Imaging: Ultrasound, CT scan, MRI, PET-CT for staging
  \item Biopsy: Core needle, excisional, or endoscopic biopsy for histopathology
  \item Immunohistochemistry to determine tumor subtype and receptor status
  \item Molecular and genetic testing (EGFR, KRAS, BRCA, MSI status)
  \item Sentinel lymph node biopsy for surgical staging
  \item Bone marrow biopsy for hematologic malignancies
  \item Genetic counseling and testing for hereditary cancer syndromes
\end{itemize}
\section{Prevention}
Avoid tobacco use and limit alcohol consumption Maintain a healthy weight through diet and exercise Eat a balanced diet rich in fruits, vegetables, and whole grains Protect skin from excessive sun exposure and avoid tanning beds Get vaccinated against cancer-causing infections (HPV, hepatitis B) Participate in recommended cancer screening programs Know your family history and consider genetic testing if indicated Avoid exposure to known carcinogens in the workplace and environment

\begin{itemize}[noitemsep]
  \item Avoid tobacco use and limit alcohol consumption
  \item Maintain a healthy weight through diet and exercise
  \item Eat a balanced diet rich in fruits, vegetables, and whole grains
  \item Protect skin from excessive sun exposure and avoid tanning beds
  \item Get vaccinated against cancer-causing infections (HPV, hepatitis B)
  \item Participate in recommended cancer screening programs
  \item Know your family history and consider genetic testing if indicated
  \item Avoid exposure to known carcinogens in the workplace and environment
\end{itemize}
\section{Treatment Options}
Surgery: Wide local excision with clear margins, lymph node dissection Radiation therapy: External beam or brachytherapy for localized disease Chemotherapy: Systemic cytotoxic drugs targeting rapidly dividing cells Targeted therapy: Drugs targeting specific molecular alterations Immunotherapy: Checkpoint inhibitors, CAR-T cells, cancer vaccines Hormone therapy: For hormone-sensitive cancers (breast, prostate) Combination approaches: Concurrent chemoradiation, sequential therapy Palliative care: Symptom management, pain control, quality of life support Clinical trials: Access to experimental therapies Surveillance: Active monitoring after definitive treatment

\begin{itemize}[noitemsep]
  \item Surgery: Wide local excision with clear margins, lymph node dissection
  \item Radiation therapy: External beam or brachytherapy for localized disease
  \item Chemotherapy: Systemic cytotoxic drugs targeting rapidly dividing cells
  \item Targeted therapy: Drugs targeting specific molecular alterations
  \item Immunotherapy: Checkpoint inhibitors, CAR-T cells, cancer vaccines
  \item Hormone therapy: For hormone-sensitive cancers (breast, prostate)
  \item Combination approaches: Concurrent chemoradiation, sequential therapy
  \item Palliative care: Symptom management, pain control, quality of life support
  \item Clinical trials: Access to experimental therapies
  \item Surveillance: Active monitoring after definitive treatment
\end{itemize}
\section{Living with It}
Regular follow-up appointments with your oncology team Manage treatment side effects with supportive medications Maintain adequate nutrition with help from a dietitian Stay physically active within your energy limits Join cancer support groups for emotional support and shared experiences Seek counseling or mental health support for anxiety and depression Communicate openly with your healthcare team about symptoms Plan for work and financial considerations during treatment Consider complementary therapies (acupuncture, massage, meditation) Build a strong support network of family and friends

\begin{itemize}[noitemsep]
  \item Regular follow-up appointments with your oncology team
  \item Manage treatment side effects with supportive medications
  \item Maintain adequate nutrition with help from a dietitian
  \item Stay physically active within your energy limits
  \item Join cancer support groups for emotional support and shared experiences
  \item Seek counseling or mental health support for anxiety and depression
  \item Communicate openly with your healthcare team about symptoms
  \item Plan for work and financial considerations during treatment
  \item Consider complementary therapies (acupuncture, massage, meditation)
  \item Build a strong support network of family and friends
\end{itemize}
\section{Outlook}
Prognosis depends on cancer type, stage at diagnosis, molecular features, and response to treatment. Early-stage cancers have significantly better outcomes, with many achieving long-term remission or cure. Survival rates have improved substantially with advances in screening, targeted therapy, and immunotherapy. Regular surveillance after treatment is essential for detecting recurrence early. Palliative care continues to advance, improving quality of life even for advanced disease.
